%% fig_ota_current.tex
%% Konventioneller OTA-Ansatz: Backend -> mehrere ECUs direkt
\begin{tikzpicture}[
  ecustyle/.style={
    draw=autosargray, fill=lightgray!50,
    rounded corners=3pt,
    minimum width=2.6cm, minimum height=0.8cm,
    align=center, font=\small
  },
  backend/.style={
    draw=autosarorange, fill=autosarorange!20,
    rounded corners=4pt,
    minimum width=3.0cm, minimum height=1.2cm,
    align=center, font=\small\bfseries
  },
  arr/.style={-{Stealth[length=5pt]}, thick}
]

\node[backend] (be) at (0,0) {Backend-Server\\{\tiny(OTA-Verwaltung)}};

\node[ecustyle] (bcu)   at (-4.5,-3) {Bremsen-SG\\{\tiny ASIL-D}};
\node[ecustyle] (radar) at (-1.5,-3) {Radar-SG\\{\tiny ASIL-B}};
\node[ecustyle] (door)  at ( 1.5,-3) {Türsteuerungs-SG\\{\tiny QM}};
\node[ecustyle] (ecu4)  at ( 4.5,-3) {Kombi-SG\\{\tiny QM}};

% Verbindungen
\draw[arr, autosarorange!70] (be.south) -- (bcu.north)
  node[midway, left, font=\tiny, xshift=-2pt] {TLS/HTTPS};
\draw[arr, autosarorange!70] (be.south) -- (radar.north)
  node[midway, left=1pt, font=\tiny] {TLS/HTTPS};
\draw[arr, autosarorange!70] (be.south) -- (door.north)
  node[midway, right=1pt, font=\tiny] {TLS/HTTPS};
\draw[arr, autosarorange!70] (be.south) -- (ecu4.north)
  node[midway, right=2pt, font=\tiny] {TLS/HTTPS};

% Problembeschriftung
\node[draw=red!70, fill=red!10, rounded corners=3pt,
      font=\tiny, text width=2.6cm, text centered]
  at (0,-4.8) {\textbf{Probleme:}\\
    Redundante Verbindungen\\Keine Synchronisierung\\
    Gro{\ss}e Angriffsfläche};

\end{tikzpicture}
